\subsection{连续线性映射的转置}

在本节中，我们假设 \((\mathrm{F},\mathrm{G})\) 和 \((\mathrm{F}_1,\mathrm{G}_1)\) 是两对成对偶的向量空间。

命题 5. --- 设 u 是 F 到 \(F_{1}\) 中的一个线性映射。下列性质等价：

\begin{enumerate}
\def\labelenumi{\alph{enumi})}
\item
  \(u\) 对弱拓扑 \(\sigma (\mathrm{F},\mathrm{G})\) 和 \(\sigma (\mathrm{F}_1,\mathrm{G}_1)\) 连续；
\item
  存在一个映射 \(v: \mathbf{G}_1 \to \mathbf{G}\) 使得
\end{enumerate}

\[
\langle u (y), z _ {1} \rangle = \langle y, v (z _ {1}) \rangle\tag{1}
\]

对所有 \(y \in \mathbf{F}\) 和 \(z \in \mathbf{G}_1\) 成立。

若这些性质成立且 \(\mathbf{F}\) 与 \(\mathbf{G}\) 之间的对偶在 \(\mathbf{G}\) 中分离，则满足 (1) 的映射 \(v\) 是唯一的，且 \(v\) 是线性的。

若 u 对弱拓扑连续，则对所有 \(z_{1} \in G_{1}\) ，F 上的线性型 \(y \mapsto \langle u(y), z_{1} \rangle\) 对 \(\sigma(F, G)\) 连续，于是 (II, 第 43 页, prop. 3) 可写成 \(y \mapsto \langle y, v(z_{1}) \rangle\) ，其中 \(v(z_{1}) \in G\) ，这表明 a) 蕴含 b)。反之，若 b) 成立，则对所有 \(z_{1} \in G_{1}\) ，线性型

\[
y \mapsto \langle y, v (z _ {1}) \rangle = \langle u (y), z _ {1} \rangle
\]

对 \(\sigma(\mathrm{F},\mathrm{G})\) 连续：由弱拓扑的定义推出 \(u\) 对 \(\sigma(\mathrm{F},\mathrm{G})\) 和 \(\sigma(\mathrm{F}_1,\mathrm{G}_1)\) 连续 (I, 第 10 页, cor. 1)。\(v\) 的唯一性由 \((\mathrm{D}_{\mathrm{II}})\) 得出，而这个唯一性蕴含 \(v\) 是线性的。

注记. --- 假设 F 与 G 之间的对偶在 G 中分离，且 \(F_{1}\) 与 \(G_{1}\) 之间的对偶在 \(G_{1}\) 中分离。若我们把 G 和 \(G_{1}\) 分别等同于 \(F^{*}\) 和 \(F_{1}^{*}\) 的子空间，则条件 a) 和 b) 等价于 \(^{t}u(\mathbf{G}_{1}) \subset \mathbf{G}\) ；v 是 u 的转置 \(^{t}u\) (A, II, §2.5) 在 \(G_{1}\) 上的限制。

我们简单地（当没有混淆可能时）称 v 为 u 的转置（关于一方面 F 与 G 之间、另一方面 \(F_{1}\) 与 \(G_{1}\) 之间的对偶），并且仍用 \(^{t}u\) 记它。

推论. --- 假设 F 与 G 之间的对偶在 G 中分离。若 u 是 F 到 \(F_{1}\) 中的一个线性映射，它对 \(\sigma(F, G)\) 和 \(\sigma(F_{1}, G_{1})\) 连续，则它的转置是 \(\mathbf{G}_1\) 到 \(\mathbf{G}\) 中的一个线性映射，且对 \(\sigma (\mathbf{G}_1,\mathbf{F}_1)\) 和 \(\sigma (\mathbf{G},\mathbf{F})\) 连续。此外，若 \(\mathbf{F}_1\) 与 \(\mathbf{G}_1\) 之间的对偶在 \(\mathbf{F}_1\) 中分离，则 \(^{t}(^{t} u) = u\) 。

在 prop. 5 中把 F 和 \(F_{1}\) 分别换成 G 和 \(G_{1}\) 即可。

命题 6. --- 假设 F 与 G（相应地 \(F_{1}\) 与 \(G_{1}\) ）之间的对偶在 G（相应地 \(F_{1}\) ）中分离。设 u 是 F 到 \(F_{1}\) 中的一个线性映射，它对 \(\sigma(F, G)\) 和 \(\sigma(F_{1}, G_{1})\) 连续。设 A 是 F 中的一个集，\(A_{1}\) 是 \(F_{1}\) 中的一个集；则：

\begin{enumerate}
\def\labelenumi{(\roman{enumi})}
\item
  我们有 \((u(\mathbf{A}))^{\circ} = ({}^{t} u)^{-1} (\mathbf{A}^{\circ})\) 。
\item
  我们有 \(\overline{{}^{t}u(A_1^\circ)}\subset (u^{-1}(A_1))^{\circ}\) ；此外，若 \(A_1\) 是闭的（对 \(\sigma(F_1,G_1)\) ）、凸的且包含原点，则我们有 \(\overline{{}^{t}u(A_1^\circ)} = (u^{-1}(A_1))^{\circ}\) 。
\end{enumerate}

设 \(z_{1} \in G_{1}\) ，关系 \(z_{1} \in (u(A))^{\circ}\) 等价于 \(\langle u(y), z_{1} \rangle \geqslant -1\) （对所有 \(y \in A\) ），而关系 \({}^{t}u(z_{1}) \in A^{\circ}\) 等价于 \(\langle y, {}^{t}u(z_{1}) \rangle \geqslant -1\) （对所有 \(y \in A\) ），于是断言 (i) 由 (1) 得出。其次，交换 \(u\) 与 \({}^{t}u\) 并把 (i) 应用于集 \(A_{1}^{\circ}\) （\(G_{1}\) 中的），我们得到

\[
\left(^ {t} u \left(\mathrm{A} _ {1} ^ {\circ}\right)\right) ^ {\circ} = u ^ {- 1} \left(\mathrm{A} _ {1} ^ {\circ \circ}\right) \supset u ^ {- 1} \left(\mathrm{A} _ {1}\right)\tag{2}
\]

由此，取极集得

\[
\left(^ {t} u \left(\mathrm{A} _ {1} ^ {\circ}\right)\right) ^ {\circ \circ} \subset \left(u ^ {- 1} \left(\mathrm{A} _ {1}\right)\right) ^ {\circ}.
\]

由双极定理 (II, 第 44 页, th. 1)，我们有 \(\overline{(^t u(A_1^\circ))} \subset (^{t}u(A_1^\circ))^{\circ \circ}\) ；最后的断言由 (2) 及双极定理得出，因为那时 \(A_1^{\circ \circ} = A_1\) 且 \(^t u(A_1^\circ)\) 是凸的并包含原点。

推论 1. --- 沿用 prop. 6 的记号，关系 \(u(\mathbf{A}) \subset \mathbf{A}_1\) 蕴含 \(^t u(\mathbf{A}_1^\circ) \subset \mathbf{A}^\circ\) ；此外，若 \(\mathbf{A}_1\) 是凸的、闭的（对 \(\sigma(\mathbf{F}_1, \mathbf{G}_1)\) ）且包含原点，则这两个关系等价。

事实上，关系 \(u(\mathbf{A}) \subset \mathbf{A}_1\) 等价于 \(\mathbf{A} \subset u^{-1}(\mathbf{A}_1)\) ，因此蕴含

\[
{ } ^ { t } u ( \mathrm{A} _ { 1 } ^ { \circ } ) \subset \overline { { { ^ { t } u ( \mathrm{A} _ { 1 } ^ { \circ } ) } } } \subset ( u ^ { - 1 } ( \mathrm{A} _ { 1 } ) ) ^ { \circ } \subset \mathrm{A} ^ { \circ }
\]

而反过来，关系 \({}^t u(\mathbf{A}_1^\circ) \subset \mathbf{A}^\circ\) 蕴含

\[
\mathrm{A} ^ {\circ \circ} \subset \left(^ {t} u (\mathrm{A} _ {1} ^ {\circ})\right) ^ {\circ} = u ^ {- 1} (\mathrm{A} _ {1} ^ {\circ \circ})
\]

后者由 (2) 得出。当 \(\mathrm{A}_1 = \mathrm{A}_1^{\circ \circ}\) 时，我们推出 \(\mathrm{A} \subset u^{-1}(\mathrm{A}_1)\) 。

推论 2. --- 设 u 是 F 到 \(F_{1}\) 中的一个线性映射，它对 \(\sigma(F, G)\) 和 \(\sigma(F_{1}, G_{1})\) 连续。则我们有

\[
\operatorname{Ker} (^ {t} u) = (\operatorname{Im} (u)) ^ {\circ},\tag{3}
\]

\[
\overline {{\operatorname{Im} (^ {t} u)}} = (\operatorname{Ker} (u)) ^ {\circ}.\tag{4}
\]

假设 F 与 G 之间以及 \(F_{1}\) 与 \(G_{1}\) 之间的对偶都是分离的；则 \(u(F)\) 在 \(F_{1}\) 中（对 \(\sigma(F_{1}, G_{1})\) ）稠密，必须且只需 \(^{t}u\) 是单射。

以 A = F 和 \(A_{1} = \{0\}\) 应用 prop. 6，并利用弱拓扑 \(\sigma(G, F)\) 和 \(\sigma(F_{1}, G_{1})\) 是 Hausdorff 的这一事实。最后的断言由 (4) 得出，其中交换 u 与 \(^{t}u\) 。
% semantic-fragments: legacy-input-seam
\relax{}
